\section{Open Questions}

The following five questions are truly open after a re-read of the paper's accessible metadata (abstract + EuropePMC record + Belov 2015 companion). None are answerable from the current Pass-2 artifacts alone; each has a concrete next step.

\subsection*{Q1. G1 vs early-S kinetic divergence under the paper's own constants}
\textbf{Basis.} The paper title explicitly advertises ``G1 and early S'' but the Pass-2 replication ran only the G1 configuration (HR suppressed). Whether the two phases produce distinguishable $t_{1/2}$ envelopes under identical damage input is untested. Without the paper's Table 1 we cannot even confirm that the paper resolves separate G1-vs-S rate constants rather than a single rate constant set with an HR on/off switch.

\textbf{Next steps.} (1) Extend \texttt{taleei\_nikjoo\_2013\_repass.py} with an early-S compartment set that adds a resection sink (rate $k_{res}$) and an HR ligation channel. (2) Sweep $k_{res}\in[0.0,0.5]$\,h$^{-1}$ and compare G1-only vs S-active total-DSB $t_{1/2}$ against Karlsson 2004 CHO / U2OS G1-vs-S synchronised foci datasets. (3) If the two envelopes overlap within experimental scatter under identical constants, the paper's G1/S distinction is a scope claim, not a kinetic claim; if they diverge $>2\times$, we need the paper's actual Table 1 to know which constants change.

\subsection*{Q2. Sensitivity of Artemis-KO 24 h residual to residual DNA-PKcs backup end-processing}
\textbf{Basis.} C7c is the only failed pass-gate: model plateaus at 30\% at 24\,h, Riballo 2004 CJ179 data sit at $\sim$18\%. The discrepancy is mechanistically explainable (Artemis-null cells retain some end-processing via DNA-PKcs autophosphorylation, per Meek et al.\ 2008), but Pass 2 explicitly did NOT silently re-fit. A principled sensitivity sweep on $k_{proc,c,backup}$ would tell us whether the paper's model can be rescued at a plausible biochemical rate or whether the gap is structural.

\textbf{Next steps.} (1) Add $k_{proc,c,backup}$ parameter, sweep in $[0.001,0.1]$\,h$^{-1}$, find the value that makes $\chi^2/dof<3$ against Riballo 2004 CJ179. (2) Compare that fitted value against DNA-PKcs autophosphorylation kinetic estimates from Meek et al.\ 2008 / Neal \& Meek 2011. (3) If the fitted value is biophysically plausible ($\sim$0.01\,h$^{-1}$), the paper's model is architecturally correct and just needs one more parameter; if not, the two-rate NHEJ topology is fundamentally under-parameterised for Artemis-deficient cells.

\subsection*{Q3. Model behaviour at ATM/ATR-mediated checkpoint arrest}
\textbf{Basis.} The paper is titled a G1/S \emph{boundary} model but Pass 2 (and, from what we can extract, the paper itself) uses pure mass-action ODE with no ATM/ATR/CHK1/CHK2 checkpoint feedback onto repair-kinase availability. If a cell arrests at the checkpoint with unrepaired complex DSBs, real biology says DNA-PKcs and Artemis abundance are upregulated (Shiloh 2003); our model would predict identical kinetics irrespective of arrest state.

\textbf{Next steps.} (1) Add a simple ATM-driven multiplier on $k_{proc,c}$ and $k_{lig,c}$ that ramps up when unrepaired complex DSB count exceeds a threshold (say 5 DSBs). (2) Rerun the C7 fits to see whether checkpoint feedback improves or worsens agreement with Beucher/Kuhne/Riballo data. (3) Read Shibata et al.\ 2011 (\textit{Radiat.\ Res.}) on ATM-dependent slow repair kinetics; if their kinetic signature matches our checkpoint-on version, the paper's mass-action-only model is incomplete for cells that actually reach the boundary.

\subsection*{Q4. Cross-species conservation of NHEJ rate constants}
\textbf{Basis.} The digitised experimental data used for C7 spans human (Beucher, Riballo CJ179), CHO (Kuhne 2000), and mouse (Asaithamby 2008 implicit) --- three species with known differences in Ku70/80 abundance, DNA-PKcs autophosphorylation kinetics, and Artemis expression. Belov 2015 uses one Table A.1 for all fits; whether Taleei \& Nikjoo 2013 do the same or resolve species-specific constants is unknown from the abstract alone.

\textbf{Next steps.} (1) Interlibrary-loan the paper PDF or contact the Karolinska group for Table 1 (only path that closes the constants-provenance gap). (2) Separately fit our 12-compartment ODE to each species dataset independently and compare fitted $(k_{proc,c}, k_{lig,c}, k_{syn})$ triples; if the three triples differ by $>30\%$ the assumption of universal constants is wrong and both replications inherit that error. (3) Cross-reference against Cucinotta 2008 track-structure model species-specific fits for a sanity floor.

\subsection*{Q5. Extension to cancer cells with checkpoint defects}
\textbf{Basis.} The paper is a mechanistic normal-cell G1/S model; LUCID-100 is a radiation-biology corpus with strong overrepresentation of cancer-radiotherapy applications. The model's construction (mass-action with fixed constants, no checkpoint feedback) makes strong implicit assumptions that break in cancer cells (p53-null, ATM-null, BRCA1-hypomorph). Testing this is a natural extension question and is not addressed by the paper.

\textbf{Next steps.} (1) Compile IMR90 (normal) vs HCT116-p53KO vs A-T fibroblast (ATM-null) residual-foci timecourses from published literature (Loebrich et al.\ 2010, Shibata et al.\ 2011). (2) Attempt to fit the same 12-compartment ODE with $(k_{proc,c}, k_{lig,c}, complex\_fraction)$ as free parameters per cell line; report whether the fitted triple lands in a biophysically plausible region for the cancer lines. (3) If it does not, the biochemical model needs a genotype-conditional layer (e.g., ATM-null bumps $complex\_fraction$ because H2AX phosphorylation is impaired) --- a natural Pass-3 extension.
